% Generado automáticamente por generar_graficos.py: no editar a mano.
\begin{table}[H]
\centering
\caption{Distribución de estudiantes según dispositivo principal y tiempo de uso nocturno antes de dormir}
\label{tab:c1}
\scriptsize
\setlength{\tabcolsep}{4pt}
\begin{tabular}{lccccccc}
\toprule
\textbf{Dispositivo principal} & \textbf{$<$30 min} & \textbf{30--59 min} & \textbf{60--119 min} & \textbf{120--179 min} & \textbf{$\ge$180 min} & \textbf{Total ($n_i$)} & \textbf{\%} \\
\midrule
Teléfono inteligente & 3 & 6 & 8 & 1 & 1 & 19 & 67,86 \\
Televisor / Smart TV & 1 & 0 & 1 & 1 & 1 & 4 & 14,29 \\
Computadora de escritorio & 1 & 0 & 0 & 1 & 1 & 3 & 10,71 \\
Computadora portátil & 0 & 1 & 1 & 0 & 0 & 2 & 7,14 \\
\midrule
\textbf{Total} & \textbf{5} & \textbf{7} & \textbf{10} & \textbf{3} & \textbf{3} & \textbf{28} & \textbf{100,00} \\
\bottomrule
\end{tabular}
\par\smallskip\parbox{0.92\linewidth}{\footnotesize\textit{Nota:} ningún informante seleccionó tableta ni la opción \guillemotleft Otro\guillemotright; ambas filas se omiten.}
\fuente
\end{table}

\begin{table}[H]
\centering
\caption{Distribución de estudiantes según tiempo de uso de pantallas antes de dormir}
\label{tab:c2}
\small
\begin{tabular}{lcccc}
\toprule
\textbf{Tiempo de uso} & \textbf{$n_i$} & \textbf{$h_i$ (\%)} & \textbf{$F_i$} & \textbf{$H_i$ (\%)} \\
\midrule
$<$ 30 min & 5 & 17,86 & 5 & 17,86 \\
30--59 min & 7 & 25,00 & 12 & 42,86 \\
60--119 min & 10 & 35,71 & 22 & 78,57 \\
120--179 min & 3 & 10,71 & 25 & 89,29 \\
$\ge$ 180 min & 3 & 10,71 & 28 & 100,00 \\
\midrule
\textbf{Total} & \textbf{28} & \textbf{100,00} &  &  \\
\bottomrule
\end{tabular}
\par\smallskip\parbox{0.92\linewidth}{\footnotesize\textit{Nota:} por tramos, 12 estudiantes (42,86 \%) usan pantallas menos de 60 min, 10 (35,71 \%) entre 60 y 119 min y 6 (21,43 \%) 120 min o más; 16 estudiantes (57,14 \%) superan los 60 min.}
\fuente
\end{table}

\begin{table}[H]
\centering
\caption{Distribución de estudiantes según duración calculada del sueño en días lectivos y sexo}
\label{tab:c3}
\small
\begin{tabular}{lcccc}
\toprule
\textbf{Duración del sueño (lun.--vie.)} & \textbf{Femenino} & \textbf{Masculino} & \textbf{Total ($n_i$)} & \textbf{\%} \\
\midrule
Menos de 5 h (muy corta) & 2 & 2 & 4 & 14,29 \\
De 5 a 6,9 h (corta) & 4 & 7 & 11 & 39,29 \\
De 7 a 9 h (óptima) & 6 & 6 & 12 & 42,86 \\
Más de 9 h (prolongada) & 1 & 0 & 1 & 3,57 \\
\midrule
\textbf{Total} & \textbf{13} & \textbf{15} & \textbf{28} & \textbf{100,00} \\
\bottomrule
\end{tabular}
\par\smallskip\parbox{0.92\linewidth}{\footnotesize\textit{Nota:} 15 estudiantes (53,57 \%) duermen menos de 7 h en días lectivos.}
\fuente
\end{table}

\begin{table}[H]
\centering
\caption{Estadísticos descriptivos de las horas calculadas de sueño según régimen semanal}
\label{tab:c4}
\footnotesize
\begin{tabular}{lccccccc}
\toprule
\textbf{Régimen} & \textbf{$\bar{x}$} & \textbf{$\tilde{x}$} & \textbf{$s$} & \textbf{RIC} & \textbf{Mín.} & \textbf{Máx.} & \textbf{CV (\%)} \\
\midrule
Días lectivos (lun.--vie.) & 6,67 & 6,67 & 1,51 & 1,75 & 4,00 & 9,33 & 22,7 \\
Fines de semana (sáb.--dom.) & 8,93 & 9,00 & 1,58 & 2,00 & 6,50 & 13,00 & 17,7 \\
Promedio ponderado (5:2) & 7,32 & 7,14 & 1,18 & 1,71 & 5,18 & 9,57 & 16,2 \\
Compensación ($\Delta t$) & 2,26 & 2,00 & 2,16 & 2,31 & -1,00 & 8,50 & 95,5 \\
\bottomrule
\end{tabular}
\par\smallskip\parbox{0.92\linewidth}{\footnotesize\textit{Nota:} horas decimales; $\Delta t = t_{\text{sueño, fin}} - t_{\text{sueño, sem}}$. Un registro inconsistente de días lectivos se imputó con la cota mínima válida (4,00 h).}
\fuente
\end{table}

\begin{table}[H]
\centering
\caption{Sensación habitual al despertar según tramo de tiempo de pantalla nocturno}
\label{tab:c5}
\footnotesize
\begin{tabular}{lcccccc}
\toprule
\textbf{Tiempo de pantalla} & \textbf{Energía} & \textbf{Cansancio leve} & \textbf{Somnolencia} & \textbf{Agotamiento} & \textbf{Total} & \textbf{Puntaje medio} \\
\midrule
$<$ 60 min & 0 & 7 & 5 & 0 & 12 & 1,58 \\
60 a 119 min & 1 & 7 & 2 & 0 & 10 & 1,90 \\
120 min o más & 0 & 6 & 0 & 0 & 6 & 2,00 \\
\midrule
\textbf{Total} & \textbf{1} & \textbf{20} & \textbf{7} & \textbf{0} & \textbf{28} & \textbf{1,79} \\
\bottomrule
\end{tabular}
\par\smallskip\parbox{0.92\linewidth}{\footnotesize\textit{Nota:} puntaje medio con la codificación del procesamiento (1 = somnolencia persistente, 2 = cansancio leve, 3 = energía suficiente); valores altos indican mejor estado matutino. Ningún informante eligió \guillemotleft agotamiento severo\guillemotright. Independencia sobre la matriz $3\times 3$: $\chi^2 = 5,7100$, $gl = 4$, $p = 0,2219$, con 77,8 \% de frecuencias esperadas menores que 5.}
\fuente
\end{table}

\begin{table}[H]
\centering
\caption{Distribución de despertares nocturnos involuntarios según ciclo académico}
\label{tab:c6}
\footnotesize
\begin{tabular}{lcccccc}
\toprule
\textbf{Ciclo académico} & \textbf{Ninguna} & \textbf{1 vez} & \textbf{2 veces} & \textbf{$\ge$3 veces} & \textbf{Total} & \textbf{\%} \\
\midrule
Básico (1.\textsuperscript{o}--3.\textsuperscript{o}) & 5 & 9 & 4 & 1 & 19 & 67,86 \\
Intermedio (4.\textsuperscript{o}--6.\textsuperscript{o}) & 1 & 3 & 1 & 0 & 5 & 17,86 \\
Finales (7.\textsuperscript{o}--10.\textsuperscript{o}) & 2 & 0 & 2 & 0 & 4 & 14,29 \\
\midrule
\textbf{Total} & \textbf{8} & \textbf{12} & \textbf{7} & \textbf{1} & \textbf{28} & \textbf{100,00} \\
\bottomrule
\end{tabular}
\par\smallskip\parbox{0.92\linewidth}{\footnotesize\textit{Nota:} semestre determinado por la materia del semestre más bajo, inscrita o no (Pregunta 3).}
\fuente
\end{table}

\begin{table}[H]
\centering
\caption{Regresión lineal simple (MCO) de las horas de sueño en días lectivos sobre los minutos de pantalla}
\label{tab:c7}
\footnotesize
\begin{tabular}{lccccc}
\toprule
\textbf{Parámetro} & \textbf{Coeficiente} & \textbf{Error est.} & \textbf{$t$} & \textbf{$p$} & \textbf{IC 95\,\%} \\
\midrule
Intercepto ($\hat\beta_0$) & 7,1896 & 0,5111 & 14,0660 & $<0{,}0001$ & [6,1390;\,8,2403] \\
Minutos de pantalla ($\hat\beta_1$) & -0,0063 & 0,0051 & -1,2221 & 0,2326 & [-0,0169;\,0,0043] \\
\bottomrule
\end{tabular}
\par\smallskip\parbox{0.92\linewidth}{\footnotesize\textit{Nota:} $n = 28$; $R^2 = 0,0543$; $R^2_{\text{aj}} = 0,0180$; $F(1,\,26) = 1,4936$, $p = 0,2326$; Durbin--Watson $= 2,074$; Jarque--Bera $= 0,632$ ($p = 0,729$); Shapiro--Wilk de residuos: $W = 0,9709$, $p = 0,6049$. Minutos imputados por marca de clase.}
\fuente
\end{table}

\begin{table}[H]
\centering
\caption{Coeficientes de correlación entre los minutos de pantalla y los indicadores de sueño}
\label{tab:c8}
\small
\begin{tabular}{llcc}
\toprule
\textbf{Coeficiente} & \textbf{Variable asociada} & \textbf{Valor} & \textbf{$p$ (bilateral)} \\
\midrule
Pearson ($r$) & Horas de sueño (lun.--vie.) & -0,2331 & 0,2326 \\
Spearman ($r_s$) & Calidad percibida (1--3) & +0,4064 & 0,0318 \\
Spearman ($r_s$) & Despertares nocturnos (0--3) & -0,0011 & 0,9957 \\
\bottomrule
\end{tabular}
\par\smallskip\parbox{0.92\linewidth}{\footnotesize\textit{Nota:} $n = 28$, $gl = n - 2 = 26$. La calidad percibida usa la escala del procesamiento (1 = somnolencia persistente, 3 = energía suficiente): un coeficiente positivo indica mejor estado matutino al aumentar los minutos de pantalla.}
\fuente
\end{table}

\begin{table}[H]
\centering
\caption{Intervalos de confianza del 95\,\% para las medias de sueño}
\label{tab:c9}
\footnotesize
\begin{tabular}{lcccc}
\toprule
\textbf{Parámetro} & \textbf{$\bar{x}$} & \textbf{$s$} & \textbf{$s/\sqrt{n}$} & \textbf{IC 95\,\%} \\
\midrule
Sueño en días lectivos ($\mu_{\text{sem}}$) & 6,6696 & 1,5124 & 0,2858 & [6,08;\,7,26] \\
Sueño en fines de semana ($\mu_{\text{fin}}$) & 8,9286 & 1,5795 & 0,2985 & [8,32;\,9,54] \\
Compensación ($\mu_{\Delta t}$) & 2,2589 & 2,1575 & 0,4077 & [1,42;\,3,10] \\
\bottomrule
\end{tabular}
\par\smallskip\parbox{0.92\linewidth}{\footnotesize\textit{Nota:} $n = 28$; $t_{0{,}975;\,27} = 2,05183$. Horas decimales.}
\fuente
\end{table}

\begin{table}[H]
\centering
\caption{Resumen de los contrastes de hipótesis ($\alpha = 0{,}05$)}
\label{tab:c10}
\footnotesize
\begin{tabular}{llcccl}
\toprule
\textbf{Prueba} & \textbf{Alternativa} & \textbf{Estadístico} & \textbf{$gl$} & \textbf{$p$} & \textbf{Decisión sobre $H_0$} \\
\midrule
1. Umbral de 7 h & $H_1: \mu_{\text{sem}} < 7$ & $t = -1,1559$ & 27 & 0,1289 & No se rechaza \\
2. Jetlag social & $H_1: \mu_{\Delta t} > 0$ & $t = 5,5401$ & 27 & $3,58\times 10^{-6}$ & Se rechaza \\
3. Modelo MCO & $H_1: \beta_1 \neq 0$ & $F = 1,494$ & (1, 26) & 0,2326 & No se rechaza \\
4. Spearman & $H_1: \rho_s \neq 0$ & $r_s = +0,4064$ & 26 & 0,0318 & Se rechaza \\
\bottomrule
\end{tabular}
\par\smallskip\parbox{0.92\linewidth}{\footnotesize\textit{Nota:} prueba 2 complementada con Wilcoxon de rangos con signo: $W = 13,0$, $p_{\text{unilateral}} = 1,13\times 10^{-5}$.}
\fuente
\end{table}
